\begin{textsample}{POS Dim 2 – ac – Score 22.00 – ac/ac\_em\_21.txt}  \label{ex:f2_pos_168}
11.4.
Aproveitamento de estudos Aproveitamento de estudos é o procedimento que as escolas da rede e as unidades certificadoras adotam para compor o resultado final de um \textbf{curso} ou etapa de ensino, por meio dos resultados alcançados em \textbf{exames}, histórico escolar, \textbf{certificação} e outros meios formais.
De acordo com a Resolução 03/2018 do Ministério da Educação, para efeito de cumprimento das exigências curriculares do ensino médio, as atividades realizadas pelos estudantes em outras \textbf{instituições}, nacionais ou estrangeiras, sejam avaliadas e reconhecidas como parte da \textbf{carga} \textbf{horária} do ensino médio, tanto da formação geral básica quanto dos \textbf{itinerários} formativos.
Assim, as atividades realizadas pelos estudantes, consideradas parte da \textbf{carga} \textbf{horária} do ensino médio, podem ser aulas, \textbf{cursos}, estágios, oficinas, trabalho supervisionado, atividades de extensão, pesquisa de campo, iniciação científica, aprendizagem profissional, participação em trabalhos voluntários e demais atividades com intencionalidade pedagógica.
Dessa forma, o aproveitamento de estudos valida os saberes e conhecimentos dos estudantes obtidos por meios formais, de modo integral ou complementar.
O aproveitamento de estudos não é obrigatório, cabendo ao estudante interessado manifestar pessoalmente ou através do seu responsável legal essa opção, na \textbf{secretaria} escolar, dentro dos prazos estipulados pela escola.
Caberá ao estudante apresentar à unidade escolar a qual está vinculado, dentre outros, os documentos comprobatórios, espelho dos resultados e/ou do boletim de desempenho escolar, a fim de que seja feita a análise e o eventual aproveitamento dos estudos, nas diferentes Áreas de Conhecimento, para posterior emissão do histórico escolar e do certificado de conclusão.
Nesse sentido, o estudante interessado poderá solicitar aproveitamento de estudos, reconhecimento de saberes e validação de competências via requerimento na \textbf{secretaria} escolar.
Esta solicitação será submetida à análise e \textbf{parecer} do Conselho escolar, conforme Regimento Interno da escola, nas seguintes situações : I – estudos realizados antes do ingresso ao Ensino Médio ; II – estudos realizados concomitantemente ao Ensino Médio.
O aproveitamento de estudos efetivar -se-á com a comprovação de que o estudante foi aprovado nos componentes que forem elencados no requerimento protocolado pelo estudante.
Neste contexto, o aproveitamento de estudos será assegurado aos estudantes mediante as seguintes situações : - forem aprovados nos componentes/unidades \textbf{cursadas} ; - prosseguirem seus estudos em que esteve vinculado ou nele reingressar ; - transferências entre \textbf{instituições} de ensino ; - estiverem em \textbf{mobilidade} curricular, ou seja, em caso de mudança de \textbf{itinerário} ( Rotas ) e, por essa razão, forem regularmente matriculados em outra \textbf{instituição}, \textbf{atendidas} as \textbf{legislações} específicas e as normas da rede ; - experiência de trabalho supervisionado ou outra experiência fora do ambiente escolar ; - demonstração de práticas ; - \textbf{cursos} oferecidos por centros ou programas de aprendizagem ; - estudos realizados em \textbf{instituições} de ensino nacionais e estrangeiras ; - \textbf{cursos} realizados por meio de educação a distância ou educação presencial mediada por tecnologias.
Os estudantes terão direito ao aproveitamento de estudos dos componentes curriculares já \textbf{cursados}, com aprovação, desde que realizados com êxito e dentro do mesmo nível de ensino, ou seja, são válidas para aproveitamento, apenas componentes de \textbf{cursos} de nível médio.
No caso de solicitação deferida, o estudante será dispensado do componente ao qual requereu, porém caso a solicitação seja indeferida, o estudante terá que \textbf{cursar} o componente ao qual requereu.
Para consideração da solicitação, será preciso que o estudante apresente : - Requerimento de aproveitamento ( disponibilizado pela \textbf{secretaria} escolar ) ; - Histórico Escolar ( contendo o nome do \textbf{curso} e/ou dos componentes/unidades curriculares, com especificação do período, frequência, \textbf{carga} \textbf{horária} e notas/conceitos ) ; - Programas, ementas e conteúdos programáticos dos componentes curriculares \textbf{cursados} com aproveitamento na escola de origem ou em outras \textbf{instituições} públicas ou particulares, exemplos : \textbf{eletivas}, \textbf{cursos} de inglês, espanhol, libras, profissionalizantes ou outros que sejam equivalentes ao componente pleiteado.
Poderão ser aproveitados os estudos realizados na modalidade presencial ou a distância. 11.5. \textbf{Oferta} de ensino à distância O cenário educacional contemporâneo mostra uma forte tendência e crescente inserção dos métodos, técnicas e tecnologias de educação a distância, dessa forma a Educação a Distância é a modalidade educacional na qual a mediação didático-pedagógica nos processos de ensino e aprendizagem ocorre com a utilização de meios e tecnologias de informação e comunicação, com estudantes e professores desenvolvendo atividades educativas em lugares ou tempos diversos.
Esta definição está presente no Decreto 5.622, de 19.12.2005 ( que revoga o Decreto 2.494/98 ), que regulamenta o Art. 80 da Lei 9.394/96 ( LDB ).
Quando as novas \textbf{Diretrizes} Curriculares Nacionais para o Ensino Médio foram aprovadas pelo Conselho Nacional de Educação ( CNE ) em 2018, criou -se a possibilidade de a escola optar pelo ensino a distância em até 30 \% na \textbf{carga} \textbf{horária} em \textbf{cursos} noturnos, 20 \% para educação diurna e até 80 \% para o EJA ( Educação de Jovens e Adultos ).
Desta forma, as \textbf{instituições} interessadas em oferecer essa modalidade de ensino precisam solicitar credenciamento específico ao Conselho Estadual de Educação, resguardada a escolha pela \textbf{instituição} educacional e/ou pela rede Estadual de ensino e cumpridos os requisitos de \textbf{oferta} da rede.
Assim, por suas características, a educação a distância supõe um tipo de ensino em que o foco está no estudante e não na turma.
Esse estudante deve ser considerado como um sujeito do seu aprendizado, desenvolvendo autonomia e independência em relação ao professor, que o orienta no sentido do aprender a aprender e do aprender a fazer.
Considerando -se que a separação física entre os sujeitos é inerente à modalidade de educação à distância, destaca -se a importância dos meios de aprendizagem e dos materiais didáticos, os quais devem ser pensados e produzidos dentro das especificidades da educação à distância e da realidade do estudante para o qual está sendo elaborado.
Dessa forma, no estado do Acre na \textbf{oferta} do ensino médio noturno e na Educação de jovens e Adultos, será considerada essa possibilidade de optar pelo ensino a distância, conforme \textbf{percentual} de \textbf{carga} \textbf{horária} indicada na Resolução nº03/2018, em atendimento a modalidade e forma de \textbf{oferta} da rede.

% matched lemmas: atender, carga, certificação, cursar, curso, diretriz, eletivo, exame, horário, instituição, itinerário, legislação, mobilidade, oferta, parecer, percentual, secretaria
\end{textsample}
